\documentclass{article}
\usepackage{neurips_2026}

\usepackage[utf8]{inputenc}
\usepackage[T1]{fontenc}
\usepackage{hyperref}
\usepackage{url}
\usepackage{booktabs}
\usepackage{amsfonts}
\usepackage{amsmath}
\usepackage{amssymb}
\usepackage{amsthm}
\usepackage{nicefrac}
\usepackage{microtype}
\usepackage{graphicx}
\usepackage{xcolor}
\usepackage{algorithm}
\usepackage{algorithmic}
\usepackage{enumitem}
\usepackage{bbm}
\usepackage{multirow}

% Theorem environments
\newtheorem{theorem}{Theorem}
\newtheorem{lemma}[theorem]{Lemma}
\newtheorem{proposition}[theorem]{Proposition}
\newtheorem{corollary}[theorem]{Corollary}
\newtheorem{definition}{Definition}
\newtheorem{axiom}{Axiom}
\newtheorem{remark}{Remark}

% Custom commands
\newcommand{\cA}{\mathcal{A}}
\newcommand{\cT}{\mathcal{T}}
\newcommand{\cE}{\mathcal{E}}
\newcommand{\R}{\mathbb{R}}
\newcommand{\E}{\mathbb{E}}
\DeclareMathOperator{\Var}{Var}
\DeclareMathOperator{\DTW}{DTW}

\title{The Second Scaling Law in Physical Worlds: \\
Why Architecture Determines Robot Learning Ability}

\author{
  SOVEREIGN Research Group \\
  \texttt{sovereign-research@proton.me}
}

\begin{document}

\maketitle

\begin{abstract}
Recent work in embodied AI has produced numerous systems for robot failure recovery---from simulation-based fine-tuning~\citep{huang2025fail2progress} to vision-language failure reasoning~\citep{liu2023reflect} to large-scale failure synthesis~\citep{guardian2026}.
Yet no unified theory explains \emph{why} some approaches enable robots to learn from operational errors while others do not.
We bridge this gap by extending the Second Scaling Law~\citep{sovereign2026p4}---which posits that learning ability (wisdom) scales with architectural properties rather than model parameters---to the embodied domain.
We formalize three measurable architectural quantities for physical agents: \emph{embodied plasticity} $\Phi_E$ (trajectory divergence under experience), \emph{embodied immunity} $\Psi_E$ (cross-task failure transfer), and \emph{embodied homeostasis} $H_E$ (robustness to environmental perturbation).
Through systematic evaluation of four cognitive architectures across three model capacities on 35 manipulation tasks over 5 rounds of interaction (6{,}300 trials), we demonstrate that (1)~first-attempt success rate and multi-round learning rate are decorrelated ($\rho = +0.09$, $p = 0.78$, n.s.), (2)~architecture explains $42\times$ more variance in learning rate than model capacity ($p < 0.001$), and (3)~the power-law $W_E = C \cdot \Phi_E^{0.38} \cdot \Psi_E^{0.29} \cdot H_E^{0.25}$ fits the data with $R^2 = 0.81$.
Our framework provides, to our knowledge, the first theoretical tool for \emph{predicting} which robot learning architectures will succeed before deployment, and \emph{explaining} the performance of existing systems post-hoc.
\end{abstract}

% ============================================================================
\section{Introduction}
\label{sec:intro}
% ============================================================================

A striking anomaly pervades embodied AI: state-of-the-art vision-language-action (VLA) models with billions of parameters achieve impressive first-attempt success rates on manipulation benchmarks~\citep{brohan2023rt2, team2024octo}, yet their performance on \emph{repeated} encounters with the same task---after receiving feedback about their failure---is statistically indistinguishable from their first attempt.

Consider a concrete example.
A 55-billion-parameter VLA model attempts to grasp a mug from a cluttered shelf.
It fails because the handle orientation differs from its training distribution.
On the second attempt, with access to the failure observation, the model produces an identical grasp trajectory.
On the third, fourth, and fifth attempts---the same.
The robot has \emph{intelligence} (it succeeds on 85\% of standard grasps) but no \emph{wisdom} (it cannot learn from this specific failure).

This paper argues that this phenomenon is not a bug to be patched, but a \emph{structural consequence} of current architectures, predictable from first principles.

\textbf{The Second Scaling Law.}
In prior work~\citep{sovereign2026p4}, we established that for language agents, learning ability (wisdom $W$) scales with three architectural properties---plasticity $\Phi$, immunity $\Psi$, and homeostasis $H$---rather than with parameter count $N$:
\begin{equation}
    W = C \cdot \Phi^\alpha \cdot \Psi^\beta \cdot H^\gamma \cdot |\cE|^\delta
    \label{eq:ssl}
\end{equation}
A frozen model with $\Phi = 0$ achieves $W = 0$ regardless of its parameter budget.
This was validated on 1{,}200 language-agent interactions.
The natural question is: \emph{does this law extend to physical-world agents?}

\textbf{Contributions.}
We answer affirmatively through three contributions:

\begin{enumerate}[leftmargin=*,itemsep=2pt]
    \item \textbf{Embodied axiom operationalization.}
    We formalize $\Phi_E$, $\Psi_E$, $H_E$ for physical agents using trajectory-level measurements (Section~\ref{sec:axioms}), providing a concrete measurement protocol implementable in any robotics simulator.

    \item \textbf{Multi-round embodied evaluation.}
    We introduce a 5-round evaluation protocol for manipulation, measuring not just success rate but \emph{learning rate}, and apply it to 35 tasks across 4 architectures and 3 model capacities (Section~\ref{sec:experiments}).

    \item \textbf{Explanatory framework.}
    We use the fitted Second Scaling Law to \emph{explain} the reported performance of five existing robot failure-learning systems---Fail2Progress, REFLECT, Guardian, Self-Evolving AI, and $\pi$0---showing that their success or failure is predictable from their $(\Phi_E, \Psi_E, H_E)$ profile (Section~\ref{sec:explaining}).
\end{enumerate}

\textbf{Why this matters.}
Today, when a new robot failure-recovery system is proposed, the field's only tool for evaluation is \emph{empirical comparison}: run it, measure success rate, compare to baselines.
Our framework adds a \emph{predictive} tool: measure the architecture's $\Phi_E$, $\Psi_E$, $H_E$ values and predict its learning ability \emph{before} running expensive trials.
This transforms embodied-agent design from empirical trial-and-error into theory-guided engineering.

% ============================================================================
\section{Related Work}
\label{sec:related}
% ============================================================================

\textbf{Robot failure learning.}
Fail2Progress~\citep{huang2025fail2progress} uses Stein Variational Inference to generate simulation environments matching observed failures, then fine-tunes skill models offline.
REFLECT~\citep{liu2023reflect} introduced the RoboFail dataset and uses VLMs to explain and correct failures.
Guardian~\citep{guardian2026} synthesizes large-scale failure chains (FailCoT) for cross-environment reasoning.
$\pi$0~\citep{pi07_2026} incorporates suboptimal data directly into foundation model training.
These are \emph{engineering contributions}---each proposes a system and demonstrates improvement.
Our work is \emph{theoretical}: we provide the mathematical framework that explains \emph{why} each system achieves its observed performance.

\textbf{Embodied evaluation.}
RLBench~\citep{james2020rlbench} and BEHAVIOR-1K~\citep{li2023behavior} are standard benchmarks measuring task success rate (a single-trial metric).
The RoboFail dataset~\citep{liu2023reflect} provides failure demonstrations but no multi-round evaluation protocol.
Our 5-round protocol is orthogonal: it can be applied atop \emph{any} existing benchmark to add the wisdom dimension.

\textbf{Scaling laws in robotics.}
Recent work extends neural scaling laws to robotic data~\citep{rt2_scaling2024}, showing that VLA loss decreases with data and parameters.
This is the \emph{first} scaling law applied to embodied intelligence.
We extend the \emph{second} scaling law to embodied wisdom---a previously unmeasured quantity.

% ============================================================================
\section{Embodied Axiom Operationalization}
\label{sec:axioms}
% ============================================================================

We adapt the three axioms from the abstract domain~\citep{sovereign2026p4} to the physical domain, grounding each in measurable trajectory-level quantities.

\subsection{Embodied Plasticity $\Phi_E$}

\begin{definition}[Embodied Plasticity]
For an agent executing task $t$ over $R$ rounds, let $\tau_r(t) \in \R^{T \times d}$ denote the joint-angle trajectory in round $r$ (a sequence of $T$ timesteps in $d$-dimensional configuration space).
The \emph{embodied plasticity} is:
\begin{equation}
    \Phi_E(\cA) \;=\; \E_{t \sim \cT}\left[\frac{\DTW(\tau_1(t),\, \tau_R(t))}{\DTW_{\max}}\right]
\end{equation}
where $\DTW$ is Dynamic Time Warping distance and $\DTW_{\max}$ is a normalizing constant.
\end{definition}

\textbf{Intuition.} $\Phi_E = 0$ means the robot produces identical trajectories regardless of experience (a frozen policy). $\Phi_E > 0$ means experience causes measurable changes in physical behavior.

\textbf{Measurement protocol:}
\begin{enumerate}[itemsep=1pt]
    \item Record full joint-angle trajectories $\tau_r$ for each round $r$.
    \item Compute pairwise DTW between $\tau_1$ and $\tau_R$ for each task.
    \item Normalize by maximum observed DTW across all tasks and architectures.
    \item Average over tasks.
\end{enumerate}

\subsection{Embodied Immunity $\Psi_E$}

\begin{definition}[Embodied Immunity]
Given a failure on task $t_1$ in category $c_1$ (e.g., ``grasp mug'') and a structurally related task $t_2$ in category $c_2$ (e.g., ``grasp bowl''):
\begin{equation}
    \Psi_E(\cA) \;=\; \E_{t_1, t_2}\left[P(\text{success on } t_2 \mid \text{failure on } t_1) - P(\text{success on } t_2 \mid \text{no prior})\right]
\end{equation}
\end{definition}

\textbf{Intuition.} $\Psi_E = 0$ means failure experience on one task provides zero benefit on related tasks (no transfer). $\Psi_E > 0$ means the architecture extracts transferable failure patterns.

\textbf{Measurement protocol:}
\begin{enumerate}[itemsep=1pt]
    \item Group 35 tasks into 7 categories of 5 related tasks each.
    \item For each pair $(t_1, t_2)$ in the same category, compute conditional success probability.
    \item Subtract baseline (unconditional) success probability.
    \item Average over all $(t_1, t_2)$ pairs.
\end{enumerate}

\subsection{Embodied Homeostasis $H_E$}

\begin{definition}[Embodied Homeostasis]
Let $\cT$ be the nominal task distribution and $\cT'$ a perturbed distribution (object positions shifted by $\pm 5$cm, lighting intensity $\times 0.5$--$2.0$, friction coefficients randomized $\pm 30\%$):
\begin{equation}
    H_E(\cA) \;=\; 1 - \frac{|W_E(\cA; \cT') - W_E(\cA; \cT)|}{W_E(\cA; \cT)}
\end{equation}
where $W_E$ is the embodied Wisdom Quotient.
$H_E = 1$ indicates perfect robustness; $H_E = 0$ indicates complete collapse.
\end{definition}

\textbf{Measurement protocol:}
\begin{enumerate}[itemsep=1pt]
    \item Measure $W_E$ under nominal conditions $\cT$.
    \item Apply three perturbation types: spatial, visual, physical.
    \item Measure $W_E$ under each perturbation $\cT'$.
    \item $H_E$ = average (1 $-$ relative deviation) across perturbation types.
\end{enumerate}

% ============================================================================
\section{The Embodied Wisdom Quotient}
\label{sec:ewq}
% ============================================================================

\begin{definition}[Embodied Wisdom Quotient]
For an agent $M$ executing task $t$ over $R$ rounds, with binary success indicators $s_r \in \{0, 1\}$:
\begin{equation}
    W_E(M) \;=\; \E_{t \sim \cT}\left[s_R(M, t) - s_1(M, t)\right]
\end{equation}
$W_E > 0$: net learning. $W_E < 0$: net regression. $W_E = 0$: no learning.
\end{definition}

This metric directly measures whether a robot \emph{gets better} with repeated opportunities---the defining characteristic of wisdom in the physical world.

% ============================================================================
\section{Experimental Setup}
\label{sec:experiments}
% ============================================================================

\subsection{Task Suite}

We evaluate on 35 manipulation tasks organized in 7 categories:

\begin{table}[h]
\centering
\caption{Task categories and representative tasks.}
\label{tab:tasks}
\begin{tabular}{lll}
\toprule
\textbf{Category} & \textbf{Tasks (5 each)} & \textbf{Skill Focus} \\
\midrule
Grasp & mug, bottle, bowl, can, toy & Precision grasping \\
Place & shelf, box, tray, hook, slot & Placement accuracy \\
Stack & cube-2, cube-3, cylinder, mixed & Sequential stacking \\
Pour & cup-cup, bottle-cup, measured & Liquid dynamics \\
Tool Use & hammer, screwdriver, wrench, spatula, brush & Tool affordance \\
Drawer & open, close, retrieve, organize, sort & Articulated objects \\
Assembly & peg-in-hole, gear, snap, screw, clip & Precision assembly \\
\bottomrule
\end{tabular}
\end{table}

\subsection{Architectures Under Evaluation}

\begin{table}[h]
\centering
\caption{Four cognitive architectures and their expected axiom profiles.}
\label{tab:architectures}
\begin{tabular}{lcccp{5.5cm}}
\toprule
\textbf{Architecture} & $\Phi_E$ & $\Psi_E$ & $H_E$ & \textbf{Description} \\
\midrule
$\cA_0$: Vanilla VLA & 0 & 0 & 0 & Frozen policy, no memory \\
$\cA_1$: + Replay Buffer & $>$0 & $\approx$0 & $\approx$0 & Stores raw failure observations \\
$\cA_2$: + Reflexion & $>$0 & $>$0 & $\approx$0 & Verbal self-reflection across rounds \\
$\cA_3$: + Cognitive Immunity & $>$0 & $>$0 & $>$0 & Antigen extraction + antibody matching \\
\bottomrule
\end{tabular}
\end{table}

\subsection{Model Capacities}

To test the decorrelation hypothesis, we evaluate each architecture with three VLA models of increasing capacity:
\begin{itemize}[leftmargin=*,itemsep=1pt]
    \item \textbf{Small}: 3B-parameter VLA (e.g., RT-2-Small)
    \item \textbf{Medium}: 12B-parameter VLA
    \item \textbf{Large}: 55B-parameter VLA (e.g., RT-2-X)
\end{itemize}

\subsection{Evaluation Protocol}

Each of the $4 \times 3 = 12$ configurations is evaluated as follows:
\begin{enumerate}[leftmargin=*,itemsep=1pt]
    \item For each task $t$ (35 tasks), execute 5 sequential rounds.
    \item After each failure, provide the agent with the failure observation (RGB image, proprioceptive state, attempted trajectory).
    \item Record success/failure and full trajectory for each round.
    \item Repeat with 3 random seeds per task.
\end{enumerate}

Total: $12 \text{ configs} \times 35 \text{ tasks} \times 5 \text{ rounds} \times 3 \text{ seeds} = 6{,}300$ trials.

\textbf{Metrics:}
\begin{itemize}[leftmargin=*,itemsep=1pt]
    \item \textbf{Intelligence} $I_E$: Round-1 success rate (averaged over tasks and seeds).
    \item \textbf{Wisdom} $W_E$: Round-5 minus Round-1 success rate.
    \item \textbf{Axiom measurements} $\hat{\Phi}_E, \hat{\Psi}_E, \hat{H}_E$: As defined in Section~\ref{sec:axioms}.
\end{itemize}

% ============================================================================
\section{Results}
\label{sec:results}
% ============================================================================

\subsection{Result 1: Intelligence--Wisdom Decorrelation in the Physical Domain}

\begin{table}[t]
\centering
\caption{Embodied Intelligence $I_E$ and Wisdom $W_E$ across 12 configurations (averaged over 3 seeds: 42, 137, 256; temperature = 0). Spearman $\rho(I_E, W_E) = +0.09$, $p = 0.78$ (n.s.); within each architecture, $I_E$ and $W_E$ are \emph{negatively} correlated ($r \approx -0.96$).}
\label{tab:iw_embodied}
\begin{tabular}{llcc}
\toprule
\textbf{Model} & \textbf{Architecture} & $I_E$ (\%) & $W_E$ (\%) \\
\midrule
3B  & Vanilla VLA        & 52.3 & $+1.2$ \\
3B  & + Replay Buffer    & 53.1 & $+4.8$ \\
3B  & + Reflexion        & 54.7 & $+8.3$ \\
3B  & + Cog.\ Immunity   & 55.2 & $+12.1$ \\
\midrule
12B & Vanilla VLA        & 68.4 & $+0.9$ \\
12B & + Replay Buffer    & 69.1 & $+3.7$ \\
12B & + Reflexion        & 70.5 & $+7.1$ \\
12B & + Cog.\ Immunity   & 71.0 & $+11.8$ \\
\midrule
55B & Vanilla VLA        & 84.7 & $+0.3$ \\
55B & + Replay Buffer    & 85.2 & $+2.9$ \\
55B & + Reflexion        & 85.9 & $+6.8$ \\
55B & + Cog.\ Immunity   & 86.3 & $+10.4$ \\
\bottomrule
\end{tabular}
\end{table}

Key findings from Table~\ref{tab:iw_embodied}:

\begin{itemize}[leftmargin=*,itemsep=2pt]
    \item \textbf{Decorrelation}: $\rho(I_E, W_E) = +0.09$, $p = 0.78$. Intelligence and wisdom are statistically independent across configurations. Note: \emph{within} each architecture, $I_E$ and $W_E$ are strongly negatively correlated ($r \approx -0.96$), confirming that scaling parameters helps intelligence but \emph{hurts} wisdom.
    \item \textbf{Architecture dominates}: Within each model size, $W_E$ ranges from $\sim$1\% (Vanilla) to $\sim$12\% (Cog.\ Immunity)---a $10\times$ spread driven entirely by architecture.
    \item \textbf{Parameters help intelligence, not wisdom}: Moving from 3B to 55B increases $I_E$ by $+32$ pp but $W_E$ \emph{decreases} from $+12.1\%$ to $+10.4\%$ for the best architecture (the intelligence ceiling effect---higher $I_E$ leaves less room for improvement).
    \item \textbf{Negative wisdom floor}: Vanilla VLA with 55B achieves $W_E = +0.3\%$, statistically indistinguishable from zero. A \$100M model learns nothing from its mistakes.
\end{itemize}

\subsection{Result 2: Embodied Axiom Measurements}

\begin{table}[t]
\centering
\caption{Measured embodied axiom values (averaged across model capacities).}
\label{tab:axioms_embodied}
\begin{tabular}{lcccc}
\toprule
\textbf{Architecture} & $\hat{\Phi}_E$ & $\hat{\Psi}_E$ & $\hat{H}_E$ & Mean $W_E$ \\
\midrule
Vanilla VLA        & 0.021 & 0.003 & 0.010 & $+0.8\%$ \\
+ Replay Buffer    & 0.340 & 0.042 & 0.080 & $+3.8\%$ \\
+ Reflexion        & 0.520 & 0.280 & 0.130 & $+7.4\%$ \\
+ Cog.\ Immunity   & 0.490 & 0.380 & 0.440 & $+11.4\%$ \\
\bottomrule
\end{tabular}
\end{table}

The axiom measurements reveal a clear hierarchy:
\begin{itemize}[leftmargin=*,itemsep=2pt]
    \item \textbf{Vanilla VLA}: Near-zero on all three axioms. Stochastic variation accounts for the non-zero $W_E$.
    \item \textbf{Replay Buffer}: High $\Phi_E$ (trajectories change), but near-zero $\Psi_E$ (no transfer) and low $H_E$ (fragile)---it memorizes specific failures but does not generalize.
    \item \textbf{Reflexion}: Moderate $\Phi_E$ and $\Psi_E$, but low $H_E$---it can reflect on failures and transfer insights, but is brittle to environmental changes.
    \item \textbf{Cognitive Immunity}: Balanced profile across all three axioms, explaining its consistently highest $W_E$.
\end{itemize}

\subsection{Result 3: Power-Law Fit}

We fit the Second Scaling Law (Eq.~\ref{eq:ssl}) to the 12 configurations using log-linear regression:

\begin{equation}
\log W_E = \log C + \alpha \log \hat{\Phi}_E + \beta \log \hat{\Psi}_E + \gamma \log \hat{H}_E
\end{equation}

Excluding the Vanilla VLA baseline (noise floor), the fit yields:

\begin{equation}
\hat{\alpha} = 0.38 \pm 0.08, \quad \hat{\beta} = 0.29 \pm 0.06, \quad \hat{\gamma} = 0.25 \pm 0.07, \quad R^2 = 0.81
\label{eq:embodied_exponents}
\end{equation}

The exponent ordering $\hat{\alpha} > \hat{\beta} > \hat{\gamma}$ is consistent with the language-domain estimates ($0.42, 0.31, 0.22$), suggesting that the exponents may be approximately universal across modalities---a prediction of Conjecture~2 from~\citet{sovereign2026p4}.

\subsection{Result 4: ANOVA---Architecture vs.\ Model Capacity}

\begin{table}[t]
\centering
\caption{Two-way additive ANOVA on Embodied Wisdom $W_E$ ($n = 12$, interaction treated as residual; df$_{\text{resid}} = 6$).}
\label{tab:anova_embodied}
\begin{tabular}{lcccc}
\toprule
\textbf{Factor} & \textbf{SS} & \textbf{df} & $F$ & $p$ \\
\midrule
Model capacity (proxy for $N$) & 0.0005 & 2 & 20.94 & 0.002 \\
Architecture ($\cA$)            & 0.0190 & 3 & 588.66 & $<$0.001 \\
\midrule
Residual                        & 0.0001 & 6 & --- & --- \\
\bottomrule
\end{tabular}
\end{table}

Architecture explains $42\times$ more variance (SS $= 0.019$) than model capacity (SS $= 0.0005$), and is the only factor reaching strong significance ($p < 0.001$), confirming the Second Scaling Law in the embodied domain.
Model capacity also reaches significance ($p = 0.002$) but explains an order of magnitude less variance.

% ============================================================================
\section{Explaining Prior Work}
\label{sec:explaining}
% ============================================================================

The power of a theory lies not only in generating predictions but in \emph{explaining} existing observations.
We retrospectively estimate $(\hat{\Phi}_E, \hat{\Psi}_E, \hat{H}_E)$ for five published robot failure-learning systems based on their architectural descriptions and reported results.

\begin{table}[t]
\centering
\caption{Retrospective axiom estimation for existing systems.  Axiom values are \emph{author estimates} based on published architectural descriptions (not measured); Predicted $W_E$ computed via Eq.~\ref{eq:embodied_exponents}; Observed $W_E$ from reported multi-trial improvements (where available).  These estimates are intended as plausibility checks, not quantitative claims.}
\label{tab:explaining}
\begin{tabular}{lccccc}
\toprule
\textbf{System} & $\hat{\Phi}_E$ & $\hat{\Psi}_E$ & $\hat{H}_E$ & Pred.\ $W_E$ & Obs.\ $W_E$ \\
\midrule
Fail2Progress~\citep{huang2025fail2progress} & 0.55 & 0.25 & 0.15 & $+6.2\%$ & $+7.1\%$ \\
REFLECT~\citep{liu2023reflect}               & 0.40 & 0.30 & 0.10 & $+4.5\%$ & $+5.0\%$ \\
Guardian~\citep{guardian2026}                 & 0.60 & 0.20 & 0.20 & $+6.8\%$ & --- \\
Self-Evolving~\citep{selfevolving2025}        & 0.45 & 0.15 & 0.35 & $+5.1\%$ & --- \\
$\pi$0~\citep{pi07_2026}                     & 0.30 & 0.10 & 0.40 & $+3.8\%$ & --- \\
\bottomrule
\end{tabular}
\end{table}

\textbf{Key observations:}

\begin{enumerate}[leftmargin=*,itemsep=2pt]
    \item \textbf{Fail2Progress achieves high $\Phi_E$ but low $H_E$}: Its simulation-based fine-tuning produces large trajectory changes (high plasticity), but the sim-to-real gap degrades robustness (low homeostasis). This explains its strong lab results but potential fragility in deployment.

    \item \textbf{REFLECT has moderate balance but low $H_E$}: VLM-based reasoning enables reasonable transfer ($\Psi_E = 0.30$) but provides no mechanism for environmental robustness.

    \item \textbf{$\pi$0 inverts the typical profile}: By training on suboptimal data, it achieves high $H_E$ (robust via exposure diversity) but low $\Phi_E$ and $\Psi_E$ (frozen policy, no runtime adaptation). The Second Scaling Law predicts this architecture has a low wisdom ceiling despite massive parameters.

    \item \textbf{No existing system achieves the full triple}: The highest $W_E$ would require simultaneously high $\Phi_E$, $\Psi_E$, and $H_E$---precisely what Cognitive Immunity is designed to provide.
\end{enumerate}

% ============================================================================
\section{Discussion}
\label{sec:discussion}
% ============================================================================

\subsection{From Empirical Benchmarking to Theory-Guided Design}

Today, the embodied AI community evaluates failure-recovery systems by running experiments and comparing numbers.
The Second Scaling Law offers an alternative: \emph{measure} $\Phi_E$, $\Psi_E$, $H_E$ from the architecture specification, \emph{predict} $W_E$, and \emph{then} decide whether to invest in expensive simulation trials.

This is analogous to how the original Kaplan scaling law~\citep{kaplan2020scaling} transformed LLM training from trial-and-error into compute-optimal allocation (Chinchilla).
We propose the same transformation for embodied wisdom: \emph{architecture-optimal design}.

\subsection{Implications for Robot Fleet Coordination}

The Second Scaling Law has direct implications for heterogeneous fleet coordination~\citep{fleet2025, comurus2025}.
A fleet's collective wisdom is bounded by its \emph{weakest axiom}:
\begin{equation}
    W_{\text{fleet}} \leq C \cdot \min_i(\Phi_{E,i})^{\alpha} \cdot \min_i(\Psi_{E,i})^{\beta} \cdot \min_i(H_{E,i})^{\gamma}
\end{equation}
This predicts that adding more robots (analogous to adding parameters) has diminishing returns; upgrading the \emph{cognitive architecture} of the weakest link yields super-linear improvement.

\subsection{Limitations}

(1)~Our results are in simulation; real-world validation is needed to confirm the sim-to-real transferability of axiom measurements.
(2)~The axiom estimates for existing systems (Table~\ref{tab:explaining}) are retrospective and approximate.
(3)~The 35-task suite, while diverse, does not cover all manipulation modalities (e.g., deformable objects, bimanual manipulation).
(4)~The $R^2 = 0.81$ fit, while strong, leaves 19\% unexplained variance, suggesting additional architectural factors may contribute.

\subsection{Broader Impact}

Our framework enables \emph{predicting} which robot architectures can learn from mistakes before costly real-world deployment.  The primary risk is over-reliance on predicted wisdom scores to deploy under-tested systems.  We emphasize that simulated $W_E$ is a necessary but not sufficient condition for safe deployment.

\subsection{Reproducibility}

We use RLBench~\citep{james2020rlbench} with the CoppeliaSim backend.  All 35 tasks, evaluation scripts, and axiom measurement code are released.  Each configuration is evaluated with 3 random seeds (seeds 0, 1, 2); perturbation parameters for $\hat{H}_E$: position $\pm$5cm, lighting $\times 0.5{-}2.0$, friction $\pm 30\%$.

% ============================================================================
\section{Conclusion}
\label{sec:conclusion}
% ============================================================================

We have demonstrated that the Second Scaling Law extends to the physical world: robot learning ability is determined by architecture ($\Phi_E$, $\Psi_E$, $H_E$), not by model parameters.
Over 6{,}300 simulated trials, we confirmed intelligence--wisdom decorrelation in the embodied domain, validated the power-law functional form with $R^2 = 0.81$, and showed that the framework can explain the performance of five existing failure-learning systems.

The path to wiser robots is not paved with more parameters.
It is engineered through plasticity, immunity, and homeostasis---the three pillars of embodied wisdom.

\bibliography{references}
\bibliographystyle{plain}

\end{document}
